% =====================================================================
% §4 Research Gap and Approach 
% =====================================================================

\section{Research Gap and Approach}
\label{sec:gap}

Section~\ref{subsec:attackers_mtd} showed the field's own surveys naming
attacker under-modelling as a primary limitation, and
Section~\ref{subsec:threat_model_eval} bore this out across a recent
five-paper cross-section, its attacker models clustering at parametric or
scripted fidelity with none reaching the procedural or behavioural rungs
at which MTD claims to operate. Seen from the attacker's side, the same
gap sharpens: where the surveyed works grant the attacker any rationality,
it is a defender-computed heuristic such as RoA~\cite{brown2023}, over
which the attacker optimises exploits it cannot sequence, adapt, or
remember---choosing the computationally best move without modelling
whether a real adversary could, or would, make it. The parametric rung, in
other words, is rationality without capability.

The means to close it already exist in MTD-adjacent literatures, in two
strands. Attack profiling
(Section~\ref{subsec:cti_profiling})~\cite{zhangY2025, rahman2025,
rodriguez2024} recovers TTP-level behaviour from raw CTI---the durable,
apex-of-the-pyramid layer the surveyed attacker models never reach,
grounded in documented APT tradecraft rather than a CVSS-derived
heuristic. This supplies the behavioural content: what the adversary does.

The second strand supplies the agent that enacts it. Bland et
al.~\cite{bland2020} build an attacker that observes partial state, acts
at a cost, and learns its strategy against a dynamic defender. Their
attacker runs over expert-validated Common Attack Pattern Enumeration and
Classification (CAPEC) patterns---evidence that the sequencing, adaptive
adversary the surveyed evaluations omit is constructible, if so far only
on narrow, hand-built cases. Outkin et al.~\cite{outkin2022} address the
parameters: they fit an attacker's per-step detection probabilities to
real MITRE ATT\&CK Evaluations data, where the surveyed models derive
theirs from CVSS. Profiled behaviour, an agent to enact it, parameters
grounded in real engagement data---the apparatus is in place and
validated; it has simply never been brought to MTD evaluation. That is the
contribution this review motivates.

\subsection{Approach}
\label{subsec:approach}

This review motivates evaluating MTD mechanisms against adversarial
profiles built by attack profiling over manually-curated CTI. The choice
of curated input is deliberate: automated extraction dominates the
literature but still falls short of reliable technique recovery in
realistic settings~\cite{buchel2025}, whereas the analyst-validated Attack
Flow corpus (Section~\ref{subsec:mitre}) trades coverage for the fidelity
an adversarial profile requires. The profiles supply the behaviour; a
Bland-style agent enacts it, parameterised against engagement data after
Outkin et al.---the two strands assembled into the adversary the
cross-section was found to lack, grounded in the tactics, techniques, and
procedures of documented APT operations rather than a defender-computed
heuristic over CVEs.

These profiles are operationalised in MTDSim~\cite{brown2023}, the
discrete-event simulator this work extends: it supports the shuffle,
diversity, and redundancy families, is agnostic to specific network
configurations, and lets the profiles run against the same parametric and
scripted baselines the cross-section surveyed---holding the defender at
baseline so that the attacker, not the orchestrator, is the variable under
test.

\subsection{Research Question}
\label{subsec:rq}

MTD mechanisms have not been evaluated against attackers modelled at
behavioural fidelity. The research question follows.
\begin{quote}
\textbf{RQ:} \textit{How do existing MTD mechanisms perform against behaviourally-grounded adversarial profiles derived from cyber threat intelligence?}
\end{quote}
The question is posed across two axes: the shuffle, diversity, and
redundancy families (Section~\ref{subsec:mtd_concept}), and profiles
differentiated by operational objective (Section~\ref{subsec:apt})---the
dimension along which APT campaigns vary, and which parametric models
collapse.

